\begin{abstract}
SHAP attributions for tabular models are usually computed by a randomized estimator under a fixed compute budget, and often from a single run. We ask when that run can stand in for ``the'' attribution of a prediction, and separate four notions of reliability: run-to-run repeatability, accuracy against the exact Shapley vector, resolvability of the feature ordering against the run-to-run noise, and the size of the explainer noise next to background-sample and model-refit variation. The evidence covers two stochastic datasets (German Credit, $d{=}20$; ACSEmployment, $d{=}16$), one background size (50 rows), one fit per model, one library version, and budgets from $\nicefrac14\times$ to $4\times$ the library default (50 instances and 20 runs per instance; $16\times$ and the exact target of the German Credit MLP use 20 instances), with two low-dimensional controls where the default is deterministic. With $\ell_1$ selection off, a single default-budget Kernel SHAP run recovers the exact top-3 feature set in at least $\nGerOneTopHitEstMinR$ of cases for all four models, and meets a pre-stated reliability rule in all of them; for the MLP, the only model here for which sampling is the practical option, the top-3 hit rate is $\nGerMlpOneTopHitEstR$ (German Credit, 20 instances) and $\nAcsMlpOneTopHitEstR$ (ACSEmployment). The conditions are narrower than that: with the library's default $\ell_1$ setting, a configuration pitfall of that implementation, the rule is not met on German Credit (relative $\ell_2$ error $\nGerLoneRelErrEstMinR$--$\nGerLoneRelErrEstMaxR$, a bias repeated runs do not remove); a single run reproduces the order of an adjacent feature pair with probability $\nGerPOrderMinR$--$\nGerPOrderMaxR$, but only $\nGerPairsResMinR$--$\nGerPairsResMaxR$ of adjacent pairs are separated by twice the run-to-run noise; and some instances are less repeatable than the average. At the 50-row background, redrawing the background changes the attribution vector far more than the explainer seed (variance ratio roughly $\nGerBgRatioMinR$--$\nGerBgRatioMaxR$ and $\nAcsBgRatioMinR$--$\nAcsBgRatioMaxR$), which is a comparison of different kinds of variation, not of accuracy. The results say nothing about larger $d$ or other background sizes.
\end{abstract}
